% Navrit Bhatti 100452705
% Aushna Shrestha 100459297
% Drishti Kapoor (100435109)
% Anmol kahlon (100413536)
% Ramneet Chandra (100360667)

\section{Conclusion}
This research paper examined the sustainability paradox and AI-optimization techniques used to improve energy efficiency while reducing the need for computational power, training, deployment, and updating with large AI models. This is typically defined within 6G environments. Thus, making this paradox not a theoretical concern, but a practical design flaw that current research hasn’t discovered a solution for yet. 

The proposed framework addresses the crisis by enabling lightweight models to handle continuous tasks and requires advanced models to be used selectively for high-demand scenarios: this avoids unnecessary use of computational power and preserves network performance. Future research should focus on evaluating models in 6G environments and make assessments regarding digital twin deployment before implementation within large-scale models. Overall, sustainable 6G design should have requirements in place to ensure constraints across both the network and AI systems are managed at a standard level. 

Literature Review: The Sustainability Paradox in AI-Driven Green Wireless Networks

Artificial intelligence (AI) is becoming a central part of next-generation wireless communication systems, especially with the development of 6G networks. It is widely seen as a tool that can improve network efficiency and reduce energy waste. However, at the same time, there is increasing concern that the energy required to train and run AI systems may cancel out some of these benefits. This creates what is often described as a sustainability paradox, where AI both supports and challenges environmental sustainability in wireless networks.

AI as a Tool for Green and Efficient Networks

A significant portion of the literature focuses on how AI can improve energy efficiency in wireless communication systems. Mao et al. (2022) explain that machine learning techniques can optimize power control, resource allocation, and network operations, helping reduce unnecessary energy consumption. Instead of relying on fixed network configurations, AI allows systems to adapt dynamically to real-time demand, which improves efficiency.

Similarly, Cui et al. (2025) emphasize that AI is expected to become a core component of 6G network architecture rather than just an optional tool. They highlight its ability to support traffic prediction, intelligent resource allocation, and energy-aware decision-making. These capabilities allow networks to operate more efficiently by avoiding overuse of resources when demand is low. Edge-based and decentralized approaches also support this trend. Fernando and Lăzăroiu (2024) show that using AI at the edge in Industrial Internet of Things (IIoT) systems can reduce communication energy costs by processing data locally instead of sending everything to centralized servers. In the same direction, Yousefpour et al. (2023) propose Green Federated Learning, where AI models are trained across distributed devices without centralizing data, reducing communication overhead and improving system efficiency. Overall, these studies suggest a clear trend: AI is being positioned as a key enabler of greener and more efficient wireless networks, especially through smarter resource management and decentralized computing.
